\documentclass[10pt]{article}

\usepackage[utf8]{inputenc}

\usepackage[T1]{fontenc}

\usepackage{amsmath}

\usepackage{amsfonts}

\usepackage{amssymb}

\usepackage[version=4]{mhchem}

\usepackage{stmaryrd}

\usepackage{bbold}

\usepackage{graphicx}

\usepackage[export]{adjustbox}

\graphicspath{ {./images/} }



\begin{document}

С. и. обладает свойством линейности как относительно интегрируемой, так и относительно интегрирующей функции (при условии существования каждого из С. и. в правой части):



\[

\begin{aligned}

& \int_{a}^{b}\left[\alpha f_{1}(x)+\beta f_{2}(x)\right] d u(x)= \\

= & \alpha \int_{a}^{b} f_{1}(x) d u(x)+\beta \int_{a}^{b} f_{2}(x) d u(x) . \\

& \int_{a}^{b} f(x) d\left[\alpha u_{1}(x)+\beta u_{2}(x)\right]= \\

= & \alpha \int_{a}^{b} f(x) d u_{1} x+\beta \int_{a}^{b} f(x) d u_{2}(x) .

\end{aligned}

\]



Свойством аддитивности С. и., вообще говоря, не обладает: из существования обоих интегралов $\int_{a}^{c} f(x) d u(x)$ и $\int_{c}^{b} f(x) d u(x)$ не следует существование интеграла $\int_{b}^{a} f(x) d u(x)$ (обратное утверждение при $a<c<b$ вместе с тем справедливо).



Если $f(x)$ ограничена на $[a, b]$, так что $m \leqslant f(x) \leqslant M$, a $u(x)$ монотонно возрастает на $[a, b]$, то найдется $\mu$, удовлетворяющее неравенствам $m \leqslant \mu \leqslant M$, такое, что для С. и. справедлива ф о р м у л а с р е д н е го Зн а ч е ни я



\[

\int_{a}^{b} f(x) d u(x)=\mu[u(b)-u(a)] .

\]



В частности, если дополнительно предположить непрерывность $f(x)$ на $[a, b]$, то найдется точка $\xi \in[a, b]$ такая, что $\mu=f(\xi)$.



С. и. $\int_{a}^{b} f(x) d u(x)$, где $u(x)$ имеет ограниченную вариацию, дает общий вид линейного непрерывного функционала $F(f)$ на пространстве непрерывных на $[a, b]$ функций (т e o p е м a $\mathrm{P}$ и с с a).



В случае; когда функция $u(x)$ имеет ограниченную вариацию, значение С. и. совпадает со значением соответствующего Лебега - Cтилтьеса интеграла.



Jum.: [1] s t i e l t j e s T h. J., Recherches sur les fractions continues, Р., 1894; [2] C м и p н о в в. И., Курс выстей матеСтильтьеса, М.- Ј’., 1936. -



СТИЛТЬЕСА ПРЕОБРАЗОВАНИЕ - интегральное преобразование вида



\[

F(x)=\int_{0}^{\infty} \frac{f(x)}{x+t} d t

\]



С. п. возникает при птерировании Лапласа преобразован̈ия и является частным случаем преобразования свертки.



Одна из формул обращения: если функция $f(t) \sqrt{t}$ непрерывна и ограничена на $(0, \infty)$, то



\[

\lim _{n \rightarrow \infty} \frac{(-1)^{n}}{2 \pi}\left(\frac{e}{n}\right)^{2 n}\left[x^{2 n} F^{(n)}(x)\right]^{(n)}=f(x)

\]



на $x \in(0, \infty)$.



C. п. о б о б ш е н н о е имеет вид



\[

F(x)=\int_{0}^{\infty} f(t) \frac{d t}{(x+t)^{\rho}}

\]



где $\rho$ - комплексное число.



C. п. ин т е г p и р о в ан но е имеет вид:



\[

F(x)=\int_{+0}^{\infty} K(x, t) f(t) d t

\]



где



\[

K(x, t)= \begin{cases}\frac{\ln \frac{x}{t}}{x-t}, & t \neq x \\ 1 / x, & t=x\end{cases}

\]



C. П. введены и для обобпенных функций. Преобразование (*) рассмотрено T. Стилтьесом (T. Stieltjes, 1894-95).



Jum.: [1] W i d d e r D. V., The Laplace transform, N. R.-L., 1946; [2] B o a s R. P., W i d d e r D. V., "Trans. Amer. Math. Soc.», 1939, v. 45, р. 1-72; [3] т и т ч м а p шI E.; Введение в теорию интегралов Фурье, пер. с англ., М.-Л.,

$1948 ;$ [4] Б р ы ч к о в Ю. А., П р у д н и к о в А. П., Инте$\begin{aligned} \text { гральные преобразования } & \text { обобщенных Функций, М., } 1977 . \\ & \text { Ю. А. Брычков, А. П. Прудникв. }\end{aligned}$



СТИНРОДА АЛГЕБРА - градуированная алгебра $A_{p}$ над полем $\mathbb{Z} p$ стационарных когомологических операций $\bmod p$. Для любого пространства (спектра пространств) $X$ группа $H^{*}(X ; \mathbb{Z} p)$ является модулем над С. a. $A_{p}$.



C. a. мультипликативно порождается $C$ тинрода операциями. Так, С. а. $A_{2}$ является градуированной ассоциативной алгеброй, мультипликативно порожденной символами $S q^{i}$ и $\operatorname{deg} S q^{i}=i$, подчиненными с оот но шения м А д е ма:



\[

S q^{a} S q^{b}=\sum_{t}\left(\begin{array}{c}

b-t-1 \\

a-2 t

\end{array}\right) S q^{a+b-t} S q^{t}

\]



$a<2 b$, так что аддитивный базис (над $\left.\mathbb{Z}_{2}\right)$ C. а. $A_{2}$ состоит из операций $S q^{i_{1}}, \ldots, S q i_{r}, i_{k} \geqslant 2 i_{k+1}$ (т. н.



\includegraphics[max width=\textwidth, center]{2023_07_09_39232d8e5387d1d5dbd1g-1}

верно и для $A_{p}$ с $p>2$. Далее,



\[

\left(A_{p}\right)^{i} \cong H^{i+n}\left(K\left(\mathbb{Z}_{p}, n\right) ; \mathbb{Z}_{p}\right), n \gg i,

\]



где $K(\mathbb{Z} p, n)$ - Эйленберга - Маклейна пространство. Умножение



\[

K\left(\mathbb{Z}_{p}, m\right) \wedge K\left(\mathbb{Z}_{p}, n\right) \rightarrow K\left(\mathbb{Z}_{p}, m+n\right)

\]



задает в С. а. диагональ $\Delta: A_{p} \rightarrow A_{p} \otimes A_{p}$, являющуюся гомоморфизмом алгебр и, следовательно, преврапающую $A_{p}$ в $X$ опфа алгебру.



Лит.: [1] С т и н р о д Н., Э п с т е йн Д., Когомологические операции, пер. с англ., М., 1983; [2] M i l n o r J., "Ann. Math.», 1958 , v. 67 , p. $150-71 ;$ [3] M о ш e p P., Т a н г о-



\begin{center}

\includegraphics[max width=\textwidth]{2023_07_09_39232d8e5387d1d5dbd1g-1(1)}

\end{center}



СТИНРОДА ДВОЙСТЕННОСТЬ - изоморфизм $p-$ мерных гомологий компактного падмножества $A$ сферы $S^{n}(n-p-1)$-мерным когомологиям дополнения (гомологии и когомологии в размерности нуль - приведенные). Рассмотрена Н. Стинродом [1]. В случае когда $A$ - открытый или замкнутый подполиәдр, аналогичный изоморфизм известен как Александера двойственность, а для любого открытого подмножества $A$ как Понтрягина двойственность. Изоморфиизм



\[

H_{p}^{c}(A ; G)=H^{n-p-1}\left(S^{n} \backslash A ; G\right)

\]



имеет место и для произвольного подмножества $A$ (д в о й с т в е н н о с т ь С и т н и к о в а); здесь $H_{p}^{c}-$ гомологии с компактными носителями Стинрода Ситникова, а $H^{q}$ - когомологии Александрова Чеха. Двойственность Александера - Понтрягина Стинрода - Ситникова - простое следствие двойственности Пуанкаре - Лефпеца и точной последовательности пары. Она справедлива не только для $S^{n}$, но и для любого многообразия, ацикличного в размерностях $p$ и $p+1$.



Jum: [1] St e e n r od N., "Ann. Math.", 1940, v. 41, p. 833-51; [2] С п т н и к о в К. А., "Докл. АН СССP", 1951 , т. 81, с. 359-62; [3] С к л я p е н к о Е. Г., "Успехи матем.

наук, 1979, т. З4, в. 6 с. $90-118 ;$ [4] М а с и У., Теория гомологий и когомологий, пер. с англ., М., 1981. СТИНРОДА ЗАДАЧА - задача реализации циклов сингулярными многообразиями; поставлена Н. Стинродом (N. Steenrod, см. [1]). Пусть $M$ - замкнутое ориентированное многообразие (топологическое, кусочно линейное, гладкое и т. Д.), и пусть $[M] \in H_{n}(M)-$ его ориентация (здесь $H_{n}(M)$ - $n$-мерная гомологий групnа многообразия $M$ ). Любое непрерывное отображение $f: M \rightarrow X$ задает элемент $f_{*}[M] \in H_{n}(X)$. С. 3. состоит в описании тех гомологич. классов из $X$, называемых р ө л и з у е м ы м и, к-рые получаются





\end{document}